\newpage
\section{État de l'Art}
\label{sec:etat_art}

La prédiction de séries temporelles (\textit{time-series forecasting}) constitue la brique centrale de nombreuses applications guidées par les données, de la santé et des réseaux énergétiques jusqu'à la surveillance industrielle et les sciences de l'environnement. Dans ces contextes applicatifs, l'objectif ne réside plus uniquement dans l'obtention de prévisions précises, mais également dans la compréhension de la dynamique décisionnelle du modèle : comment parvient-il à une prédiction donnée, et quelles modifications permettraient d'atteindre une trajectoire future plus idéale ? Cette exigence d'explicabilité est d'autant plus critique dans le cadre de la prédiction de séries temporelles multivariées, où la variable cible dépend conjointement de son propre historique, de covariables externes, de décalages temporels (\textit{lags}) et d'interactions croisées entre variables.

Bien que l'avènement des modèles d'apprentissage profond ait nettement amélioré les performances prédictives \cite{petropoulos2022forecasting, kong2025deep, wang2024deep, lv2024efficient}, ces architectures demeurent de véritables « boîtes noires » difficiles à interpréter pour les experts métier.

\subsection{Explicabilité pour la prédiction de séries temporelles}

Les modèles profonds de prédiction se sont imposés en raison de leur capacité à capturer des dépendances temporelles non linéaires et des interactions complexes entre variables \cite{petropoulos2022forecasting, kong2025deep, wang2024deep}. Dans un cadre multivarié, cette puissance d'expression s'avère particulièrement pertinente dès lors que la cible dépend de plusieurs covariables observées au cours du temps. Néanmoins, cette complexité altère la lisibilité des prédictions.

Pour y remédier, la majorité des méthodes d'explicabilité développées pour les séries temporelles reposent sur l'attribution d'importance (\textit{feature attribution}) ou la saillancie (\textit{saliency}). Ces approches cherchent à mesurer la contribution relative des variables, des pas de temps ou des sous-séquences à la décision finale \cite{saadallah2021explainable, pan2021series, nguyen2025tshap, choi2025vector}. À titre d'exemple, l'approche \textit{Series Saliency} apprend un masque explicatif sur les dimensions temporelles et thématiques afin de conjuguer prédiction et interprétabilité \cite{pan2021series}. De même, des adaptations de SHAP ont été formulées spécifiquement pour le traitement des séries temporelles en classification, régression et prédiction \cite{nguyen2025tshap, choi2025vector}.

Si ces méthodes d'attribution sont précieuses pour identifier les zones critiques sur lesquelles le modèle concentre son attention, elles souffrent d'un manque d'actionnabilité. Elles n'indiquent pas à l'utilisateur comment le signal d'entrée (historique ou variables exogènes) devrait être altéré pour obtenir un comportement prédictif souhaité. Cette limite constitue la motivation première de l'étude des explications contrefactuelles.

\subsection{Explications contrefactuelles pour la prédiction}

Une explication contrefactuelle cherche l'altération minimale à appliquer aux données d'entrée afin de conduire le modèle vers une sortie cible prédéfinie. Dans la littérature générale de l'IA explicable (XAI), l'évaluation de ces explications s'appuie habituellement sur des critères tels que la validité, la proximité, la parcimonie (\textit{sparsity}), l'actionnabilité et la plausibilité \cite{guidotti2024counterfactual, moreira2025benchmarking}. Appliqués aux séries temporelles, ces critères doivent intégrer une contrainte supplémentaire liée à la structure temporelle : la modification contrefactuelle doit non seulement rester proche de l'entrée originale, mais également préserver la cohérence temporelle et inter-variables.

Toutefois, la formulation d'explications contrefactuelles en prédiction s'avère nettement plus complexe qu'en classification. En classification, le résultat visé correspond à une classe discrète. En prédiction, la sortie est continue et s'étend souvent sur un horizon multi-pas : la cible souhaitée peut alors prendre la forme d'une trajectoire exacte, d'un intervalle d'acceptabilité ou du résultat d'une intervention ciblée. La littérature existante présente une grande diversité d'approches :
\begin{itemize}
    \item \textbf{Approches par bornes et intervalles :} \textit{ForecastCF} \cite{wang2023forecastcf} est l'un des premiers travaux à avoir formalisé l'explication contrefactuelle pour la prédiction en définissant la cible sous forme de bornes inférieure et supérieure sur l'horizon de prévision. \textit{COMET} \cite{wang2024comet} étend ce principe au cadre multivarié (pour la prédiction du glucose) en y intégrant des contraintes du domaine métier (écrêtage, limites d'activité et contraintes sur l'historique), démontrant ainsi que la validité seule ne suffit pas si l'explication n'est pas réaliste sur le plan applicatif.
    \item \textbf{Approches par trajectoire cible :} À l'inverse, des méthodes comme \textit{CET-X} \cite{kinjo2025cetx} visent une trajectoire cible exacte en modifiant prioritairement les variables exogènes pour forcer la prédiction à s'en rapprocher.
    \item \textbf{Approches fondées sur des valeurs extrêmes ou la recherche heuristique :} \textit{DiffusionCF} \cite{deng2024unraveling} traite la prédiction de valeurs extrêmes en s'appuyant sur la théorie des valeurs extrêmes pour définir une cible contrefactuelle non critique. Enfin, d'autres travaux combinent causalité au sens de Granger, régression quantile et algorithmes génétiques pour générer des scénarios futurs plausibles \cite{zuin2024navigating}.
\end{itemize}

\subsection{Méthodes contrefactuelles connexes et inférence causale}

Au-delà de la prédiction multi-pas, d'autres approches contrefactuelles ont été développées pour des tâches connexes sur séries temporelles. \textit{CounTS} \cite{yan2023self} propose un modèle auto-interprétable reposant sur l'inférence contrefactuelle (abduction, action, prédiction). \textit{NTD-CFE} \cite{sun2025no} s'appuie sur l'apprentissage par renforcement sans accès au jeu d'entraînement (\textit{model-agnostic}). D'autres travaux adressent spécifiquement la classification de séries temporelles \cite{delaney2021instance}, la génération de contrefactuels rares/épars \cite{lang2023sparse} ou la détection d'anomalies \cite{srinivasan2024ae}. Bien qu'apportant des concepts utiles en matière d'actionnabilité et de sobriété, ces méthodes ne sont pas directement configurées pour évaluer des modèles de prédiction multivariés multi-horizons.

Il convient également de distinguer ces méthodes d'explication des travaux issus de l'inférence causale et du contrôle synthétique (tels que TGV-CRN, CDA, MVDLM ou la modélisation sous contraintes d'injection d'eau) \cite{fujii2025estimating, min2024domain, li2024compositional, wen2025dynamical}. Ces derniers s'attachent à estimer l'effet réel d'une intervention physique ou d'une politique de traitement au cours du temps. À l'opposé, les explications contrefactuelles \textit{post-hoc} cherchent à expliquer le fonctionnement d'un modèle d'IA déjà entraîné en identifiant les ajustements d'entrée nécessaires pour satisfaire une contrainte prédictive donnée.

\subsection{Synthèse et positionnement}

En résumé, la littérature sur les contrefactuels appliqués à la prédiction reste morcelée. Les méthodes diffèrent non seulement par leurs stratégies d'optimisation, mais aussi par la définition de la tâche, la façon dont la prédiction souhaitée est spécifiée (trajectoire exacte vs bornes), les variables autorisées à la modification, et les métriques d'évaluation retenues. Ce manque de standardisation rend la comparaison directe complexe et souligne le besoin de benchmarks unifiés pour évaluer le comportement des différentes familles d'explications contrefactuelles face à un même problème de prédiction multivariée.

\subsection{Positionnement du cas hydrogéologique}

Le cas d'usage traité ici se place à l'intersection de plusieurs des familles qui précèdent, et s'en écarte sur deux points qu'il faut expliciter avant d'aborder le protocole d'évaluation.

Il relève d'abord des approches par bornes plutôt que par trajectoire exacte. La question posée par un gestionnaire de la ressource n'est pas \emph{quelle chronique exacte de pluie aurait produit ce niveau}, mais \emph{quel régime climatique aurait suffi à ramener la nappe au-dessus du seuil d'alerte}. La cible naturelle est donc un domaine d'acceptabilité sur l'horizon de prévision, formulation qui rejoint celle de \textit{ForecastCF} \cite{wang2023forecastcf}, et le seuil se lit sur l'échelle standardisée à sept classes plutôt que sur une valeur absolue de niveau, laquelle n'est pas comparable d'un ouvrage à l'autre.

Il impose ensuite une contrainte de plausibilité plus forte que la proximité géométrique usuelle, pour la raison exposée en section~\ref{sec:xai} : les variables climatiques d'entrée sont liées par des relations physiques qu'une perturbation libre viole aussitôt. Cette exigence est de même nature que celle qui a conduit \textit{COMET} \cite{wang2024comet} à intégrer des limites de domaine métier, et elle sépare nettement les deux réponses implémentées ici : la paramétrisation à faible dimension garantit la plausibilité par construction, là où la recherche d'un analogue observé la garantit par la provenance, dans la lignée des approches par instances \cite{delaney2021instance, ates2021comte}. La littérature ne tranche pas entre ces deux garanties, et c'est l'un des objets de l'évaluation proposée en section~\ref{sec:methodo}.

Il se distingue enfin des travaux d'inférence causale évoqués ci-dessus par la nature de la question posée. On ne cherche pas à estimer l'effet réel d'une intervention sur l'aquifère, ce qui relèverait de la modélisation hydrogéologique et non de la XAI, mais à caractériser ce sur quoi un modèle prédictif appris fonde sa sortie. La disponibilité, dans la plateforme, d'un modèle mécanistique calibré indépendamment sur la même station offre à cet égard une position peu commune dans la littérature XAI : un second modèle, interprétable par construction, contre lequel confronter le contrefactuel produit sur le modèle appris.
